\section{Conclusiones}
\label{sec:conclusiones}

Este trabajo formuló el problema de priorizar, gestionar y derivar
solicitudes de atención y lo resolvió con una estructura de datos
justificada para cada necesidad: colas enlazadas por nivel y un min-heap
por plazo como representaciones de la cola de pendientes, un árbol AVL como
índice por identificador, una pila para deshacer, Heap-Sort para los
reportes, un árbol de decisión para clasificar la urgencia y un grafo de
sedes con BFS y Dijkstra para la derivación. Registrar, atender, buscar,
cancelar y deshacer cuestan $O(\log n)$ en el peor caso.

La implementación cubre los doce requerimientos funcionales y pasa 69
verificaciones automáticas. Los experimentos mostraron que el AVL evita la
degeneración del BST con identificadores consecutivos, con una altura de 14
frente a 19\,999 y un tiempo 280 veces menor para 20\,000 solicitudes, y que
la política por plazo elimina la inanición del nivel menos urgente bajo
sobrecarga, a cambio de un costo por operación mayor que el de la cola por
niveles.

Como trabajo futuro se plantea: cargar el protocolo de triage y la red de
sedes desde archivos; permitir cambiar la urgencia de una solicitud con una
operación de disminución de clave sobre un heap indexado; migrar las
pruebas a un marco como JUnit; y evaluar el sistema con datos de llegada
más realistas.
